\documentclass{article}
\usepackage[final]{neurips_2026}
\makeatletter
\renewcommand{\@noticestring}{RoboSynChallenge 2026}
\makeatother
\usepackage[T1]{fontenc}
\usepackage{booktabs}
\usepackage{graphicx}
\usepackage{float}
\usepackage{amsmath}
\usepackage[hidelinks]{hyperref}
\title{SZU-VCC Synapse: Adapting $\pi_{0.5}$\\for Bimanual Manipulation}
\author{Jiahao Li \\
Shenzhen University, SZU-VCC Synapse \\
\texttt{ljh285915@gmail.com}}
\begin{document}
\maketitle

\begin{abstract}
We present our $\pi_{0.5}$ approach to the ten RoboSynChallenge finals tasks.
Each task is first learned from synthetic demonstrations and then adapted
to the released physical demonstrations. Synthetic data cover changes in
object placement, geometry, appearance, and contact dynamics. Physical
adaptation uses trajectory-level training and validation splits, with
checkpoint selection guided by validation loss and generated-action error.
We report complete 150-episode simulation evaluations for items handover,
click bell, and handle basket. Their 10,000-update synthetic checkpoints
achieve 87.3\%, 62.7\%, and 36.7\% success, respectively. We also describe
the data preparation, an image-guided assembly refinement module, and the
shared inference interface used to deploy the task-specific finals weights.
\end{abstract}

\section{Overview}
RoboSynChallenge studies how synthetic manipulation data can support
generalization to physical robots with limited demonstrations
\citep{robosynchallenge}. Our finals system uses OpenPI $\pi_{0.5}$
\citep{pi05,openpi} for all ten tasks. We train a separate policy for each
task so that its synthetic data, physical adaptation, instruction, and
normalization statistics can be chosen together.

The training pipeline has two stages. Synthetic demonstrations first teach
the policy the task sequence and expose it to variations in the scene.
We then initialize from a selected synthetic checkpoint and fine-tune
using the released physical demonstrations. Most tasks retain synthetic
samples during this stage to preserve task coverage while adapting to
the physical images and joint trajectories. A shared deployment entry
loads the selected policy for each task and converts its predictions into
joint-space commands.

\section{Synthetic Data and Policy Training}
\subsection{Demonstration generation}
For most tasks, the training pool contains 640 accepted demonstrations:
400 base trajectories, 55 position variations, 41 size variations, 96
visual variations, 28 combined variations, and 20 dynamics variations.
The physical ranges are designed separately for each task. Position and
size changes alter reach and grasp geometry; visual changes vary surfaces,
lighting, and distractors; dynamics changes vary mass and friction.
Combined scenes bring several changes into the same episode.

Demonstrations are generated by executing the collection controller in
simulation. Each accepted trajectory must complete the task and pass
checks for finite state/action values, synchronized camera streams,
frame alignment, and control cadence. We inspect three-camera videos
before scaling up collection. Later data versions also randomize the
colors and textures of the three visible walls. The dataset records
the background version of each trajectory so that the actual visual
coverage can be traced.

Handle-basket collection additionally checks joint motion at the recorded
10\,Hz cadence. An audit identified 18 trajectories requiring a slower
withdrawal transition. We re-executed those trajectories and retained
the other 622, keeping the original group quotas and scene seeds.
The resulting 640-trajectory pool contains 115,283 frames.

Item assembly uses 719 curated simulation sources, combining approach
and grasp segments with recovery demonstrations. Segment boundaries
are checked to retain useful motion while excluding incorrect connection
tails. This gives the policy additional examples of approaching and
aligning the two parts.

\subsection{Policy and optimization}
We initialize from OpenPI $\pi_{0.5}$ and fully fine-tune the model on
each task's synthetic data. Training uses the flow-matching objective.
The model receives the high camera, left-wrist camera, right-wrist
camera, a task instruction, and proprioception. Images follow the
OpenPI 224-by-224 preprocessing path; the instruction and discretized
state use the model's tokenizer with a 200-token limit.

The active robot state has 14 components: six left-arm joints, the left
gripper, six right-arm joints, and the right gripper. Joint targets are
represented relative to the current joint state during training, while
gripper targets remain absolute. The action vector is padded to the
model's internal dimension of 32. At inference, inverse normalization
and the current joint state recover absolute joint targets.

\subsection{Image-guided assembly refinement}
We also developed an optional refinement module for the final connection
motion in item assembly. The $\pi_{0.5}$ policy performs the approach and
grasp. Wrist-camera geometry estimates track the relative positions and
axes of the two tube ends, with temporal filtering to stabilize the
estimates. A handoff gate activates near the alignment region. The
controller then converts the endpoint and orientation error into small,
bounded joint corrections through forward and inverse kinematics,
refreshing the visual estimate after each correction. The module is
packaged separately with its geometry model and kinematic configuration.

\section{Simulation Evaluation}
\subsection{Protocol}
We completed the same evaluation structure for items handover, click bell,
and handle basket: 60 diagnostic scenes, 50 combined scenes, and 40 stress
scenes per checkpoint. Diagnostic scenes isolate sources of variation;
combined scenes include regular, boundary, and extrapolated combinations;
stress scenes target task-specific shifts such as unseen appearance,
challenging geometry, and contact conditions. Scene seeds are separated
from the training pool. The 5,000- and 10,000-update checkpoints use the
same frozen scenes and execution protocol within each task.

Evaluation runs at 10\,Hz control and 100\,Hz physics. The policy predicts
50 actions and executes ten before receiving a new observation. An
episode ends on the task's native success condition or at 1,000 control
steps. Every complete episode remains in the result table. These local
development evaluations are used to compare synthetic checkpoints before
physical-data adaptation.

\begin{table}[t]
\centering
\caption{Complete local simulation evaluation of simulation-trained $\pi_{0.5}$ checkpoints. Each row contains 150 episodes. Diagnostic, combined, and stress columns show successful episodes over the full split.}
\label{tab:simulation}
\small
\begin{tabular}{lrrrrrr}
\toprule
Task & Update & Diagnostic & Combined & Stress & Total & Success \\
\midrule
Items handover & 5,000 & 1/60 & 4/50 & 0/40 & 5/150 & 3.3\% \\
Items handover & 10,000 & 54/60 & 40/50 & 37/40 & 131/150 & 87.3\% \\
Click bell & 5,000 & 7/60 & 15/50 & 19/40 & 41/150 & 27.3\% \\
Click bell & 10,000 & 42/60 & 32/50 & 20/40 & 94/150 & 62.7\% \\
Handle basket & 5,000 & 14/60 & 10/50 & 13/40 & 37/150 & 24.7\% \\
Handle basket & 10,000 & 26/60 & 20/50 & 9/40 & 55/150 & 36.7\% \\
\bottomrule
\end{tabular}
\end{table}


\subsection{Results}
The 10,000-update checkpoint improves overall success in all three tasks.
Items handover reaches 131/150 successes, including 54/60 diagnostic,
40/50 combined, and 37/40 stress scenes. Click bell improves from 41/150
to 94/150; its largest gain is in the diagnostic split. Handle basket
improves from 37/150 to 55/150. Its stress result decreases from 13/40
to 9/40 even as the diagnostic and combined results improve, showing that
longer training does not improve every type of scene equally.

Time constraints limited the extent of simulation evaluation across the
remaining tasks. We used small-scale tests for development and focus the
quantitative results in this report on the three complete 150-episode
evaluations above.

\section{Adaptation to Physical Demonstrations}
\subsection{Data split and mixed training}
The released physical dataset contains 60 demonstrations per task.
We split complete trajectories into 48 training and 12 validation
trajectories. Splitting by trajectory keeps frames from the same
demonstration together. For drawer, the checked-cloth group forms the
validation set, while the white and blue groups are used for training.
For water pouring, one excluded training trajectory leaves 47 training
and 12 validation trajectories. Normalization statistics are computed
from the training partition.

Physical adaptation starts from the selected task-specific synthetic
checkpoint. Most runs use batches of 16, with eight synthetic and eight
physical samples. Water pouring uses physical-only batches of 16,
initialized from its 5,000-update synthetic checkpoint. The common peak
learning rate is $5\times10^{-6}$. Training runs on H800 GPUs, using one
GPU or two-device sharding according to the task configuration.

\subsection{Checkpoint selection}
We compare saved checkpoints on a fixed validation partition. The
selection diagnostics are the validation flow-matching loss and the
error between generated actions and demonstration actions. Generated
actions are compared after restoring physical units, with joint and
gripper errors tracked separately. Fixed inputs and random seeds are
used across candidate checkpoints, and trajectories receive equal
weight so that longer demonstrations do not dominate the comparison.
For receding-horizon control, action comparisons focus on the first ten
predicted steps, which are the steps executed before replanning.

Candidates with lower validation loss or lower action error under the
task's fixed protocol are preferred. The retained adaptation checkpoints
are listed in Table~\ref{tab:adaptation}. Each is distributed with the
normalization statistics used by its training run.

\begin{table}[H]
\centering
\caption{Retained physical-adaptation checkpoints. Update counts start from each task's selected synthetic initialization.}
\label{tab:adaptation}
\small
\begin{tabular}{lr@{\hspace{1cm}}lr}
\toprule
Task & Update & Task & Update \\
\midrule
Items handover & 1,500 & Click bell & 1,500 \\
Handle basket & 1,500 & Item assembly & 1,500 \\
Drawer open place & 3,000 & Water pouring & 1,500 \\
Sample loading & 1,500 & Table rearrangement & 1,500 \\
Mixer operating & 4,000 & Manipulate pipette & 2,000 \\
\bottomrule
\end{tabular}
\end{table}

\section{Deployment and Conclusion}
The ten finals task bindings use a shared $\pi_{0.5}$ entry under
\texttt{policy/szu\_synapse}. It loads the task checkpoint, verifies its
paired normalization file, and accepts three RGB camera views,
proprioception, and the task instruction. The physical recording's
32-dimensional state is mapped to the active joints and grippers using
indices 0--6 and 16--22. The entry returns 14-dimensional absolute
joint/gripper commands and executes ten actions per observation update
at the 10\,Hz control rate.

The public release contains inference parameters, normalization files,
per-file checksums, installation instructions, and task-level execution
commands. A bundled tokenizer keeps inference self-contained. Model and
code revisions identify the exact submitted artifacts.

Our approach combines structured synthetic coverage with task-specific
physical adaptation. The complete simulation results show gains from
longer synthetic training while also exposing task-dependent weaknesses
under stress. The physical demonstrations then provide the observations
and action conventions needed for the final deployment stage. Official
real-robot evaluation will measure the resulting task success.

\bibliographystyle{plainnat}
\bibliography{references}
\end{document}
